\documentclass[10pt,a4paper]{article}
\usepackage[margin=22mm]{geometry}
\usepackage{amsmath,amssymb,amsthm,booktabs,longtable,hyperref}
\newtheorem{theorem}{Theorem}\newtheorem{proposition}{Proposition}
\newcommand{\gen}[1]{\langle#1\rangle}
\setlength{\emergencystretch}{2em}
\title{Literal presentations and local periods in hexagonal extensions}
\author{An application to Fernandes--Leemans--Ivi\'c Weiss}
\date{Private mathematical draft; complete review pending}
\begin{document}\maketitle
The retained primary is the 21-page author edition of
\emph{Hexagonal extensions of toroidal maps and hypermaps} \cite{source}.
We preserve its parameter order, literal presentations and regular/chiral
conventions. The original constructions and questions are attributed to
its authors. All group orders below are full \emph{type-preserving}
automorphism orders; correlations are separate.

The complete earlier marked-kernel common paper and the complete new
\emph{Rotation certificates and affine period families} supplement are
paid premises. The latter contains the entire accepted cubic, square
and exact finite proofs formerly located in the L121 application, now
shared with this application. No sibling application proof is free.
The standard torus and faithful affine Coxeter conventions remain
explicit classical inputs. We do not reprove the prior classifications
quoted as source Tables 1--2 or Theorems 3.1--3.3.

\section{One certificate for each literal finite row}
For each parameter row let $P=\gen{x,y,z\mid\mathcal R}$ be exactly the
source's rotation presentation. The entire six presentation builders
occur in the required program \path{presentation_engine.g}.
Parameters are read from \path{table_roster.json}, not reconstructed
from a claimed order. Put $a_0=1,a_1=x,a_2=y,a_3=z$ and
$P_i=\gen{a_j^{-1}a_k:j,k\ne i}$.
The engine establishes an exact isomorphism from the abstract
presentation to its permutation group. This is an upper-bound
computation as well as a finite image. In T8.8, complete coset
enumeration over $\gen{y,z}$ gives index 17280; the abstract local
presentation has order 150. The resulting permutation image has
order 2592000, equal to 17280\,$\cdot$\,150, so the upper and lower
bounds coincide. ACE is bounded by 500000 cosets,5000000 workspace
words and 60 seconds. The first general attempt was stopped at 120
seconds; its receipt is not treated as evidence against finiteness.

For every row the engine checks all 256 intersections of the 16
pair-difference subgroups, all 16 normalized flag counts, generation
for each of the four rank-three residues and the six actual rank-two
residue incidence graphs. The shared theorem gives thinness, residual
connectivity and precisely two rotation-group chamber orbits.
If simultaneous inversion satisfies every literal relator it is an
involutory automorphism. The faithful two-copy semidirect completion
is checked for all 256 marked intersections and 16 coset-Helly identities.
Otherwise a specified failed inverse relator certifies chirality and
that $P$ itself is the full type-preserving group.

\begin{theorem}[The computed literal rows]
All 66 finite rows below define finite regular or chiral hypertopes,
with the exact full orders and actual marked local groups displayed.
Each has the source's asserted regular/chiral status. The full orders
in source Table 5 rows 8 and 10 must be doubled for the literal
presentations in this retained edition.
\end{theorem}
\begin{proof}
Apply the shared rotation and local-kernel propositions to the exact
per-row outputs. The abstract upper bound and generator maps are
separate from the subgroup and flag tests. Each local group is
identified by an onto map from its complete two-generator presentation
with matching order; its full translation kernel is also computed.
The regular completion has order $2|P|$ and the chiral full group has
order $|P|$. No assertion is transferred between rows merely because
their numerical orders agree. The required software and entire row
roster are explicit computational premises.
\end{proof}
R and C mean regular and chiral. The four local orders are rotation
orders in omitted-type order 0,1,2,3. Seven inherited-family rows are
handled separately. Printed structural descriptions are audited separately below. They are
never identified merely from an order.
\small\begin{longtable}{llrrcl}
\toprule Row&Literal parameters&Printed $g$&Proved $g$&Type&Local rotation orders\\\midrule\endhead
T3.1 & $1,2;\,1,2$ & 1,344 & 1,344 & C & 12/42/12/42 \\
T3.2 & $1,2;\,2,1$ & 20,160 & 20,160 & C & 12/42/12/42 \\
T3.3 & $2,0;\,1,2$ & 672 & 672 & C & 12/42/12/24 \\
T3.4 & $2,0;\,1,3$ & 4,368 & 4,368 & C & 12/78/12/24 \\
T3.5 & $2,0;\,1,4$ & 16,128 & 16,128 & C & 12/126/12/24 \\
T3.6 & $2,0;\,2,3$ & 13,680 & 13,680 & C & 12/114/12/24 \\
T4.1 & $2,0;\,2,0$ & 240 & 240 & R & 6/24/6/24 \\
T4.2 & $2,0;\,6,0$ & 720 & 720 & R & 6/72/6/24 \\
T4.3 & $3,0;\,3,0$ & 2,916 & 2,916 & R & 6/54/6/54 \\
T4.4 & $3,0;\,4,0$ & 241,920 & 241,920 & R & 6/96/6/54 \\
T4.5 & $3,0;\,1,1$ & 324 & 324 & R & 6/18/6/54 \\
T4.6 & $3,0;\,2,2$ & 41,472 & 41,472 & R & 6/72/6/54 \\
T4.7 & $1,1;\,1,1$ & 108 & 108 & R & 6/18/6/18 \\
T4.8 & $2,2;\,2,2$ & 13,271,040 & 13,271,040 & R & 6/72/6/72 \\
T4.9 & $3,0;\,1,3$ & 33,696 & 33,696 & C & 6/78/6/54 \\
T4.10 & $4,0;\,1,2$ & 12,096 & 12,096 & C & 6/42/6/96 \\
T4.11 & $6,0;\,1,2$ & 756,000 & 756,000 & C & 6/42/6/216 \\
T4.12 & $1,2;\,3,6$ & 2,016 & 2,016 & C & 6/126/6/42 \\
T4.13 & $1,2;\,4,4$ & 36,288 & 36,288 & C & 6/288/6/42 \\
T4.14 & $1,2;\,6,6$ & 2,268,000 & 2,268,000 & C & 6/648/6/42 \\
T4.15 & $3,5;\,2,1$ & 672 & 672 & C & 6/42/6/42 \\
T4.16 & $1,4;\,2,1$ & 2,016 & 2,016 & C & 6/42/6/126 \\
T4.17 & $1,2;\,1,2$ & 672 & 672 & C & 6/42/6/42 \\
T5.1 & $2,0;\,2,0;\,2,0$ & 240 & 240 & R & 6/24/24/12 \\
T5.2 & $2,0;\,3,0;\,3,0$ & 1,296 & 1,296 & R & 6/54/54/12 \\
T5.3 & $2,0;\,4,0;\,4,0$ & 15,360 & 15,360 & R & 6/96/96/12 \\
T5.4 & $2,0;\,6,0;\,2,2$ & 2,880 & 2,880 & R & 6/72/72/12 \\
T5.5 & $3,0;\,2,0;\,2,0$ & 1,296 & 1,296 & R & 6/24/24/27 \\
T5.6 & $3,0;\,2,0;\,3,0$ & 13,824 & 13,824 & R & 6/24/54/27 \\
T5.7 & $3,0;\,2,0;\,4,0$ & 165,888 & 165,888 & R & 6/24/96/27 \\
T5.8 & $3,0;\,2,0;\,5,0$ & 2,592,000 & 5,184,000 & R & 6/24/150/27 \\
T5.9 & $3,0;\,2,0;\,2,2$ & 3,888 & 3,888 & R & 6/24/72/27 \\
T5.10 & $3,0;\,2,0;\,4,4$ & 248,832 & 497,664 & R & 6/24/288/27 \\
T5.11 & $3,0;\,3,0;\,1,1$ & 972 & 972 & R & 6/54/18/27 \\
T5.12 & $3,0;\,1,1;\,1,1$ & 324 & 324 & R & 6/18/18/27 \\
T5.13 & $3,0;\,1,1;\,3,3$ & 2,916 & 2,916 & R & 6/18/162/27 \\
T5.14 & $3,0;\,2,2;\,2,0$ & 3,888 & 3,888 & R & 6/72/24/27 \\
T5.15 & $4,0;\,2,0;\,2,0$ & 15,360 & 15,360 & R & 6/24/24/48 \\
T5.16 & $2,0;\,1,2;\,2,1$ & 336 & 336 & C & 6/42/42/12 \\
T5.17 & $2,0;\,1,3;\,2,5$ & 2,184 & 2,184 & C & 6/78/78/12 \\
T5.18 & $2,0;\,1,4;\,1,2$ & 336 & 336 & C & 6/42/42/12 \\
T5.19 & $2,0;\,1,4;\,3,6$ & 8,064 & 8,064 & C & 6/126/126/12 \\
T5.20 & $2,0;\,2,3;\,3,2$ & 6,840 & 6,840 & C & 6/114/114/12 \\
T5.21 & $3,0;\,1,2;\,1,2$ & 275,562 & 275,562 & C & 6/42/42/27 \\
T6.1 & $2,0;\,2,0$ & 384 & 384 & R & 4/24/24/12 \\
T6.6 & $3,0;\,1,1$ & 144 & 144 & R & 4/18/18/12 \\
T6.7 & $6,0;\,1,1$ & 576 & 576 & R & 4/72/18/12 \\
T6.8 & $3,0;\,1,3$ & 58,968 & 58,968 & C & 4/54/78/12 \\
T6.9 & $1,2;\,1,2$ & 2,688 & 2,688 & C & 4/42/42/12 \\
T6.10 & $1,2;\,2,1$ & 1,008 & 1,008 & C & 4/42/42/12 \\
T6.11 & $1,2;\,3,1$ & 58,968 & 58,968 & C & 4/42/78/12 \\
T7.1 & $3;\,3,0;\,1,1$ & 360 & 360 & R & 12/12/18/18 \\
T7.2 & $3;\,6,0;\,1,1$ & 23,040 & 23,040 & R & 12/12/18/72 \\
T7.3 & $3;\,1,2;\,1,1$ & 1,512 & 1,512 & C & 12/12/18/42 \\
T7.4 & $3;\,1,4;\,1,1$ & 90,720 & 90,720 & C & 12/12/18/126 \\
T7.5 & $3;\,2,4;\,1,1$ & 774,144 & 774,144 & C & 12/12/18/168 \\
T7.6 & $4;\,2,0;\,1,1$ & 4,320 & 4,320 & R & 24/24/18/24 \\
T7.7 & $4;\,1,1;\,1,1$ & 3,456 & 3,456 & R & 24/24/18/18 \\
T8.4 & $2,0;\,1,1;\,4,0;\,3,0$ & 165,888 & 165,888 & R & 96/54/18/24 \\
T8.5 & $3,0;\,2,0;\,1,1;\,2,0$ & 1,296 & 1,296 & R & 18/24/24/54 \\
T8.6 & $3,0;\,2,0;\,1,1;\,3,0$ & 13,824 & 13,824 & R & 18/54/24/54 \\
T8.7 & $3,0;\,2,0;\,1,1;\,4,0$ & 165,888 & 165,888 & R & 18/96/24/54 \\
T8.8 & $3,0;\,2,0;\,1,1;\,5,0$ & 5,184,000 & 5,184,000 & R & 18/150/24/54 \\
T8.9 & $2,0;\,1,2;\,1,1;\,6,0$ & 15,120 & 15,120 & C & 18/72/42/24 \\
T8.10 & $2,0;\,1,2;\,1,2;\,2,0$ & 352,800 & 352,800 & C & 42/24/42/24 \\
T8.11 & $2,0;\,1,2;\,1,1;\,0,3$ & 7,620,480 & 7,620,480 & C & 18/504/42/24 \\
\bottomrule\end{longtable}
\normalsize

\section{Actual local periods and corrections}
The standard torus presentations have orders $6(a^2+ab+b^2)$, or
$3(a^2+ab+b^2)$ for the triangular hypermap. Twelve ordinary period
relations become stronger local quotients in the global presentation.
For instance, T7.1 has actual local rotation order 18 where $(3,0)$
would give 54. Each exact local kernel is recorded below, in the ordered
translation words of the shared local-kernel proposition. It specifies
the marked quotient, not an unlabelled torus selected by order.
\small\begin{longtable}{lcrrl}
\toprule Row&Omitted type&Nominal local order&Actual order&Exact kernel basis\\\midrule\endhead
T4.2 & 1 & 216 & 72 & $(6,0);(2,2)$ \\
T4.12 & 1 & 378 & 126 & $(21,0);(4,1)$ \\
T4.15 & 3 & 294 & 42 & $(7,0);(2,1)$ \\
T5.4 & 1 & 216 & 72 & $(6,0);(2,2)$ \\
T5.17 & 2 & 234 & 78 & $(13,0);(3,1)$ \\
T5.18 & 1 & 126 & 42 & $(7,0);(2,1)$ \\
T5.19 & 2 & 378 & 126 & $(21,0);(4,1)$ \\
T6.6 & 1 & 54 & 18 & $(3,0);(1,1)$ \\
T6.7 & 1 & 216 & 72 & $(6,0);(2,2)$ \\
T7.1 & 3 & 54 & 18 & $(3,0);(1,1)$ \\
T7.2 & 3 & 216 & 72 & $(6,0);(2,2)$ \\
T8.9 & 1 & 216 & 72 & $(6,0);(2,2)$ \\
T8.11 & 1 & not a period & 504 & $(42,0);(8,2)$ \\
\bottomrule\end{longtable}
\normalsize

There is an additional literal issue in source Section 8, physical p.17.
Its final $h$-word is $y^{-1}(yz)^2=zyz$, whose cube is already trivial
from $y^3=z^2=1$. For T8.11, $v=(0,3)$ thus adds no period.
The literal global group nevertheless has the printed order 7620480.
Its actual local rotation order is 504, with kernel basis $(42,0),(8,2)$
and determinant 84. We certify this actual quotient from the three
triangle relations and its full lattice kernel, without inserting a
likely intended different word. All other positive Table 8 rows have $h=0$.

T5.8 has rotation order 2592000 and regular full order 5184000;
T5.10 has rotation order 248832 and regular full order 497664.
Both pass every intersection, flag and actual-residue test. The printed
full totals are 2592000 and 248832. Their structural abbreviations cannot
be asserted unchanged when their implied orders conflict; the exact
marked presentations and local models specify the corrected objects.

Adding the actual local kernel relations changes no global group,
since they already hold in it. Conversely the original period
relations hold in each exact local quotient. Thus the same global
group is universal for these \emph{actual} local presentations.
A claim with uncollapsed preassigned periods is stronger and is not
made for the listed exceptions.

\section{Shared affine families and inherited finite rows}
For Table 6 rows 2--5 and all $s\ge3$, the complete shared cubic proof
gives $E/sD$ of full order $48s^3$, its marked residues and a faithful
degree-$4s$ CPR representation. In its notation $(a,b,d,c)$ use
\[x=ba,\qquad y=bc,\qquad z=bd.\]
Then $(y^{-1}z^{-2})^2=(c(bd)^2b)^2$ is the fixed local relation,
and $(x^{-1}z^{-2})^s=(a(bd)^2b)^s$ is the variable relation.
The ordered parameter pair is $((2,0),(s,0))$. The full cubic toroid
has order $96s^3$; the common proof retains its explicit index-two
embedding and the source CPR labels. The endpoint $s=2$ is separately
certified as T6.1, not inferred from the faithful-$4s$ argument.

Table 8 row 1 is the entire shared square family with parameters
$((2,0),(1,1),(s,s),(2,0))$ and full order $1152s^4$, for $s\ge1$.
The fixed affine group, integral parabolic lattices, period kernel,
intersections and flag transitivity are proved in the paid common
proof. Table 8 rows 2--3 are the complete shared finite cases of orders 720
and 46080. Their actual first local period has the degree-three quotient
proved there, rather than the uncollapsed $(3t,0)$ type. The shared orientation-adapter proof gives exact rotation/reflection
presentation identities and explicit marker permutations alongside the
unchanged proof migration; matching orders alone is not the argument.

\section{Structural columns: verified models and explicit limits}
The following ledger preserves every one of the 73 printed group descriptions.
V records a verified isomorphism to the explicitly stated target, M an exact
shared affine or finite marked model without a claimed identification of the
printed extension action, R an unverified source extension description, and
O an order-incompatible description. An R or M entry is not an assumed group
identification and is not a premise of the geometry or universal-presentation
proof. Thus the table makes the remaining structural-column obligations visible;
this manuscript does not claim every printed structure string proved.

The required \path{structural_groups.py} constructs the same full group from
the certified rotation generators and the checked inversion automorphism.
For the 29 V rows it constructs the stated target and verifies an actual
bijective group homomorphism. It records the source generators, their images
and the target generators. The target denoted $A\mathbin{\times_{C_2}}S_3$
is the subgroup of $A\times S_3$ where the two quotient signs agree, using
$A/A'=C_2$; this is the ``fibre'' notation below. The group
$\operatorname{P\Sigma L}_2(27)$ is the preimage of the order-three subgroup
of $\operatorname{Aut}(\operatorname{PSL}_2(27))/\operatorname{PSL}_2(27)$.
For $\operatorname{P\Sigma L}_2(49)$ the program explicitly adjoins the field
map $u\mapsto u^7$ to the fractional-linear $\operatorname{PSL}_2(49)$
action. This distinguishes the field action from the inequivalent diagonal
extension $\operatorname{PGL}_2(49)$. Wreath products use their full stated
permutation action. Classical finite-group constructors and exact isomorphism
algorithms are declared software premises; no source extension action is
inferred from composition factors.

For T8.11 a smaller factor certificate replaces an expensive direct
isomorphism search. The two minimal normal subgroups have orders 2520 and
504. Their opposite centralizers have orders 5040 and 1512, commute,
intersect trivially, and generate the full group. Separate bijections identify
them with $S_7$ and $\operatorname{Aut}(\operatorname{PSL}_2(8))$.
Multiplication therefore proves the displayed direct-product identification.
The first bounded direct search timed out; no failure claim was inferred.
\small\begin{longtable}{@{}lp{.40\linewidth}p{.46\linewidth}@{}}
\toprule Row & Printed group column & Exact scope here\\\midrule\endhead
T3.1 & $2 \mathbin{.} L_{2}(7) : 2 : 2$ & R: source description retained as unverified; no extension action asserted.\\
T3.2 & $2 \mathbin{.} A_{7} : 2 : 2$ & R: source description retained as unverified; no extension action asserted.\\
T3.3 & $L_{2}(7) : 2 \times 2$ & V: PGL(2,7) x C2\\
T3.4 & $L_{2}(13) \times 2 \times 2$ & V: PSL(2,13) x C2 x C2\\
T3.5 & $Q_{8} \mathbin{.} (L_{2}(7) \times 2) : S_{3}$ & R: source description retained as unverified; no extension action asserted.\\
T3.6 & $L_{2}(19) : 2 \times 2$ & V: PGL(2,19) x C2\\
T4.1 & $A_{5} : 2 \times 2$ & V: S5 x C2\\
T4.2 & $S_{5} \times S_{3}$ & V: S5 x S3\\
T4.3 & $3^{1+2} \times 3 : S_{3} : S_{3}$ & R: source description retained as unverified; no extension action asserted.\\
T4.4 & $\operatorname{Aut}(L_{3}(4))$ & V: Aut(PSL(3,4))\\
T4.5 & $3^{1+2} : 2 : S_{3}$ & R: source description retained as unverified; no extension action asserted.\\
T4.6 & $2^{1+6} : 3^2 : S_{3} : S_{3}$ & R: source description retained as unverified; no extension action asserted.\\
T4.7 & $3^{1+2} : 2^2$ & R: source description retained as unverified; no extension action asserted.\\
T4.8 & $2^{5+6} : A_{5} \times 3 : S_{3} : S_{3}$ & R: source description retained as unverified; no extension action asserted.\\
T4.9 & $L_{3}(3) : S_{3}$ & V: Aut(PSL(3,3)) fibre\_C2 S3\\
T4.10 & $U3(3) : 2$ & V: Aut(PSU(3,3))\\
T4.11 & $U3(5) : S_{3}$ & V: Aut(PSU(3,5))\\
T4.12 & $2 \mathbin{.} L_{2}(7) : S_{3}$ & R: source description retained as unverified; no extension action asserted.\\
T4.13 & $U3(3) : S_{3}$ & V: Aut(PSU(3,3)) fibre\_C2 S3\\
T4.14 & $3 \mathbin{.} U3(5) : S_{3}$ & R: source description retained as unverified; no extension action asserted.\\
T4.15 & $2 \mathbin{.} L_{2}(7) : 2$ & R: source description retained as unverified; no extension action asserted.\\
T4.16 & $2 \mathbin{.} L_{2}(7) : S_{3}$ & R: source description retained as unverified; no extension action asserted.\\
T4.17 & $2 \mathbin{.} L_{2}(7) : 2$ & R: source description retained as unverified; no extension action asserted.\\
T5.1 & $S_{5} \times 2$ & V: S5 x C2\\
T5.2 & $S_{3} \times S_{3} \times S_{3} : S_{3}$ & V: S3 wreath S3\\
T5.3 & $2^{1+6} : (A_{5} : 2)$ & R: source description retained as unverified; no extension action asserted.\\
T5.4 & $S_{5} \times 2^2 : S_{3}$ & V: S5 x S4\\
T5.5 & $6 : S_{3} : S_{3} : S_{3}$ & R: source description retained as unverified; no extension action asserted.\\
T5.6 & $2^{1+6} : 3 : S_{3} : S_{3}$ & R: source description retained as unverified; no extension action asserted.\\
T5.7 & $2^{1+6} : 6 : S_{3} : S_{3} : S_{3}$ & R: source description retained as unverified; no extension action asserted.\\
T5.8 & $2 \mathbin{.} (A_{5} \times (A_{5} \times A_{5})) : 3 : 2$ & O: printed order incompatible; no structural identification asserted.\\
T5.9 & $S_{3} \times 3 : S_{3} : S_{3} : S_{3}$ & R: source description retained as unverified; no extension action asserted.\\
T5.10 & $2^{1+6} : (3 : S_{3} : S_{3} : 3) : S_{3}$ & O: printed order incompatible; no structural identification asserted.\\
T5.11 & $3^{1+2} : 6 : S_{3}$ & R: source description retained as unverified; no extension action asserted.\\
T5.12 & $3^2 \times S_{3} : S_{3}$ & R: source description retained as unverified; no extension action asserted.\\
T5.13 & $3^{1+2} \times 3 : S_{3} : S_{3}$ & R: source description retained as unverified; no extension action asserted.\\
T5.14 & $S_{3} \times 3 : S_{3} : S_{3} : S_{3}$ & R: source description retained as unverified; no extension action asserted.\\
T5.15 & $4 \mathbin{.} (2^4 : A_{5}) : 2 \times 2$ & R: source description retained as unverified; no extension action asserted.\\
T5.16 & $L_{2}(7) : 2$ & V: PGL(2,7)\\
T5.17 & $L_{2}(13) \times 2$ & V: PSL(2,13) x C2\\
T5.18 & $L_{2}(7) : 2$ & V: PGL(2,7)\\
T5.19 & $Q_{8} \mathbin{.} L_{2}(7) : S_{3}$ & R: source description retained as unverified; no extension action asserted.\\
T5.20 & $L_{2}(19) : 2$ & V: PGL(2,19)\\
T5.21 & $3^{2+6} : 7 : 3 : 2$ & R: source description retained as unverified; no extension action asserted.\\
T6.1 & $2^{1+4} \times 2 : S_{3}$ & R: source description retained as unverified; no extension action asserted.\\
T6.2 & $S_{3} \times S_{3} \times S_{3} : S_{3}$ & M: complete shared marked model; printed extension action not identified.\\
T6.3 & $2^{1+6} : 2^2 : S_{3}$ & M: complete shared marked model; printed extension action not identified.\\
T6.4 & $5^3 : 2 : 2 : 2 : S_{3}$ & M: complete shared marked model; printed extension action not identified.\\
T6.5 & $2^3 : S_{3} : S_{3} : S_{3} : S_{3}$ & M: complete shared marked model; printed extension action not identified.\\
T6.6 & $S_{3} \times 2^2 : S_{3}$ & V: S3 x S4\\
T6.7 & $2^4 : S_{3} : S_{3}$ & R: source description retained as unverified; no extension action asserted.\\
T6.8 & $L_{2}(27) : 3 : 2$ & V: Aut(PSL(2,27))\\
T6.9 & $2^6 : 7 : 3 : 2$ & R: source description retained as unverified; no extension action asserted.\\
T6.10 & $L_{2}(7) \times 3 : 2$ & V: PGL(2,7) x C3\\
T6.11 & $L_{2}(27) : 3 \times 2$ & V: PSigmaL(2,27) x C2\\
T7.1 & $A_{5} : S_{3}$ & V: S5 fibre\_C2 S3\\
T7.2 & $2^4 : A_{5} : 2 : 2 : S_{3}$ & R: source description retained as unverified; no extension action asserted.\\
T7.3 & $L_{2}(8) : 3$ & V: Aut(PSL(2,8))\\
T7.4 & $L_{2}(8) \times A_{5} : 3$ & V: Aut(PSL(2,8)) x A5\\
T7.5 & $2^9 \mathbin{.} L_{2}(8) : 3$ & R: source description retained as unverified; no extension action asserted.\\
T7.6 & $A_{6} : 2 \times S_{3}$ & V: S6 x S3\\
T7.7 & $2^{1+4} : 3 : S_{3} : S_{3}$ & R: source description retained as unverified; no extension action asserted.\\
T8.1 & $s^4 : 2^{1+4} : S_{3} : S_{3}$ & M: complete shared marked model; printed extension action not identified.\\
T8.2 & $A_{5} \times 2 : S_{3}$ & M: complete shared marked model; printed extension action not identified.\\
T8.3 & $2^4 : A_{5} : 2 \times 2 : 2 : S_{3}$ & M: complete shared marked model; printed extension action not identified.\\
T8.4 & $2^{1+6} : S_{3} : S_{3} : S_{3} : S_{3}$ & R: source description retained as unverified; no extension action asserted.\\
T8.5 & $S_{3} \times S_{3} \times S_{3} : S_{3}$ & V: S3 wreath S3\\
T8.6 & $2^{1+6} : 3 : S_{3} : S_{3}$ & R: source description retained as unverified; no extension action asserted.\\
T8.7 & $2^{1+6} : S_{3} : S_{3} : S_{3} : S_{3}$ & R: source description retained as unverified; no extension action asserted.\\
T8.8 & $2 \mathbin{.} (A_{5} \times (A_{5} \times A_{5})) : 2 : S_{3}$ & R: source description retained as unverified; no extension action asserted.\\
T8.9 & $3 \times S_{7}$ & V: C3 x S7\\
T8.10 & $L_{2}(49) \times 3 : 2$ & V: PSigmaL(2,49) x C3\\
T8.11 & $S_{7} \times (L_{2}(8) : 3)$ & V: S7 x Aut(PSL(2,8))\\
\bottomrule\end{longtable}
\normalsize

\section{The diagram list and its limit}
The ten patterns in source Figure 1 represent sixteen fixed edge-labelled
cases after expanding its parameter $p\in\{3,4,5,6\}$. Here is the
complete elementary reduction. A connected rank-four nonlinear diagram
has one of five underlying graphs: a star, a triangle with a pendant
edge, a four-cycle, $K_4$ minus one edge, or $K_4$. Every rank-three
triangle must be spherical or Euclidean, and at least one must be
Euclidean. The retained triangle criterion is
$1/p+1/q+1/r\ge1$, with equality for a toroidal face and omitted
edges assigned label two. An edge with label greater than six would
have a neighbouring edge of label at least three by connectedness,
violating this criterion. Thus labels two through six suffice.

A star must have arm labels $(3,3,6)$, $(3,4,4)$ or $(4,4,4)$.
A triangle with a pendant edge has all triangle labels three and
pendant label three through six. In a four-cycle, a label six forces
its neighbours to be three and permits the opposite label three
through six. With no six, an affine adjacent pair must be $(4,4)$;
labels five cannot then occur. The possibilities are an adjacent pair
of fours and two threes, three fours and one three, or four fours.
Finally, in $K_4$ minus an edge and in $K_4$, all edges belong to
triangles and hence all have label three. This gives respectively
$3+4+7+1+1=16$ cases, exactly the source's ten parametrized patterns.
The optional rational enumeration in \path{diagram_types.py} checks
this case split; it is not a premise of the written argument.
The reduction is a list of possible diagrams, not a proof that every
case has a finite universal hypertope.

\section{Full source boundaries and questions}
Section 2's coset construction and C-plus terminology retain Tits and
the earlier-literature attributions. The shared theorem supplies the
finite rotation-to-geometry direction with its actual hypotheses.
Source Tables 1--2 and Theorems 3.1--3.3 quote earlier classifications
and remain retained inputs. Experiments and suggested correlation
constructions around diagrams 4--5, and bounded failed searches in
Section 8, are documentary; they are not newly proved classifications
or nonexistence theorems. The source restricts its investigation to
the first three diagrams.

Problems 9.1--9.2 and Conjectures 9.1--9.3 are separate from positive-row
recovery. Previously reviewed Problem 9.2 work constructs finite
quotient geometries for every integer parameter and proves the literal
universal group infinite for all parameters at least six. Its exact
old proofs and reviews are retained separately; no fresh solution,
priority or current-open status is claimed. The four small universal
orders now have the T8.5--T8.8 certificates. Finite quotients do not
imply finite universal groups. None of these computations proves the
source lists exhaustive or proves its negative conjectures.

\begin{thebibliography}{9}
\bibitem{source} M.~E.~Fernandes, D.~Leemans and A.~Ivi\'c Weiss,
\emph{Hexagonal extensions of toroidal maps and hypermaps}, retained
21-page author edition (2018 institutional file).
\bibitem{shared} \emph{Rotation certificates and affine period families},
complete supplement to \emph{Supported kernels, translation lattices
and finite marked quotients}, private proof migration.
\bibitem{later} M.~E.~Fernandes, D.~Leemans, C.~A.~Piedade and
A.~Ivi\'c Weiss, \emph{Two families of locally toroidal regular
4-hypertopes arising from toroids}, twelve-page author edition.
\end{thebibliography}
\end{document}
